%% ---------------------------------------------------------------------------
%% Revised Section 2 (condensed for the 35-page limit; every citation of the
%% submitted version is kept).
%% ---------------------------------------------------------------------------
\section{Related Work}

Three lines of research meet in this paper. Explainability for language models provides the token-level attributions we evaluate and the confidence-based error detectors they must be measured against (Section~\ref{explainability_lm}). Explainability for graph neural networks provides the graph explainers and the fidelity metrics from which our four evaluation dimensions are built (Section~\ref{explainability_graphs}). Graph representations of text provide the topologies on which those explainers run and the precedent for pairing a graph network with a language model, as in our LLM-as-teacher setting (Section~\ref{text_trees}). We review each line for what it has established about explanations as signals of model behaviour and for the comparisons it leaves open.

\subsection{Explainability in Language Models} \label{explainability_lm}

Surveys of explainability for LLMs agree on the map of the field and on its open problems. Luo and Specia~\cite{luo2024understandingutilizationsurveyexplainability} divide the methods into local analyses of individual predictions and global analyses of the model, place feature attribution (perturbation-, gradient- and vector-based) at the centre of the local family, and note that perturbation methods such as LIME and SHAP are limited by their cost; they also review how explanations are used beyond understanding, for model editing, capability enhancement and controllable generation, that is, as signals that feed back into the model. Zhao et al.~\cite{zhao2024surveyllm-xai} organise the same methods by training paradigm (fine-tuning versus prompting) and by scope, with feature attribution, attention, example-based and natural-language explanations as the local categories, and list among the open challenges the absence of ground-truth explanations, so that faithfulness has to be measured with removal-based metrics such as comprehensiveness and sufficiency, for which well-established criteria are still lacking. The same gaps persist in later surveys of the area~\cite{palikhe2025survey-llm-xai}, and they shape our protocol. Without ground truth, we judge an explanation by its effect on the model's output when kept or removed, and by whether the prediction is correct. Because a query costs more on a language model than on a graph network, we fix the explainer budget in model queries rather than in time.

Benchmarks confirm these gaps: in ALMANACS~\cite{mills2025almanacssimulatabilitybenchmarklanguage} current LLM explanation methods fail simulatability tests, and in PRobELM~\cite{yuan2025probelmplausibilityrankingevaluation} plausibility rankings frequently diverge from human intuitions. Efforts to simplify gradient-based attributions into more interpretable forms~\cite{rodrigues2024gad} reflect the demand for explanation methods that produce clearer, more actionable signals.

Error detection itself has an established baseline family that any explanation-based signal must be compared with: the maximum softmax probability~\cite{hendrycks2017baseline}, the top-2 margin and the predictive entropy, together with learned confidence estimators~\cite{corbiere2019confidnet} and nearest-neighbour trust scores~\cite{jiang2018trust}. Their usefulness depends on calibration, which modern networks often lack~\cite{guo2017calibration}. Whether a distilled student inherits that confidence is open: distillation trains the student on the teacher's class probabilities~\cite{gou2021kd}, yet students often fail to match the teacher's predictive distribution~\cite{stanton2021kd}; Section~\ref{sec:calibration} measures what a graph surrogate inherits from a language-model teacher. Evaluation of attribution methods has its own pitfalls, from saliency maps that do not depend on the model~\cite{adebayo2018sanity} to deletion-based metrics that measure the retraining distribution rather than the attribution~\cite{hooker2019roar}; our protocol therefore evaluates every explanation-based detector against majority-class and confidence baselines.

\subsection{Explainability for Graph Neural Networks} \label{explainability_graphs}

GNN explanations take the form of node, edge and node-feature importance masks or subgraphs~\cite{agarwal2023graphxai}. Perturbation-based explainers infer importance from prediction changes: SubgraphX~\cite{yuan2021explainabilitygraphneuralnetworks} searches connected subgraphs with Monte Carlo tree search, and GraphSVX~\cite{duval2021graphsvxshapleyvalueexplanations} generalises Shapley values to node coalitions. Self-explainable graph convolutional architectures with frequency-adaptive mechanisms~\cite{wei2024selfexplainable} offer intrinsic interpretability as an alternative to post-hoc methods. GraphFramEx~\cite{amara2022graphframex} and GNNX-BENCH~\cite{kosan2024gnnxbenchunravellingutilityperturbationbased} benchmark post-hoc GNN explainers systematically across graph structures, and feature perturbation augmentation~\cite{brocki2023fpa} addresses the reliability of the fidelity metrics themselves, a concern that motivates our multi-dimensional evaluation. Because the reliability of a graph model depends on the structure it is given~\cite{xue2025llm-gnn-xai}, we evaluate several text-to-graph conversions.

\subsection{Graph Representations of Text} \label{text_trees}

Graph neural networks enter NLP through the representation chosen for the text, and three lines of prior work supply the candidates. Syntactic graphs come from parsers: Universal Dependencies gives dependency structure a standard that holds across languages~\cite{de-marneffe-etal-2021-universal}, and transition-based constituency parsers now reach state-of-the-art accuracy in several languages~\cite{bauer-manning-2025-high}, so both kinds of tree can be extracted at scale. Co-occurrence graphs need no parser: TextGCN~\cite{yao2019textgcn} builds a graph from word co-occurrence within sliding windows and trains a graph convolutional network on it for text classification. A third line combines a graph network with a pre-trained language model, for instance a graph convolutional network and BERT joined through co-attention for fake news detection~\cite{zhang2024gbca}. Our topologies are drawn from the first two lines (constituency and dependency trees; window and skip-gram graphs), and the third is the closest precedent for a graph network whose node features come from a language model.
